% !TEX root = Chicken Chanko Nabe Meal Prep.tex
\documentclass[14pt]{extarticle}
\usepackage[utf8]{inputenc}
% --- REMOVE PAGE NUMBERING ---
\pagestyle{empty}
% --- MARGIN SETUP ---
\usepackage[left=0.75in, right=0.75in, top=0.5in, bottom=0.5in]{geometry}
% --- LIST SETTINGS ---
\usepackage{enumitem}
\setlist[itemize]{leftmargin=1em}
% Labels are centered in a fixed-width box so plain bullets line up with the
% wider prep icons (\veg, \measure) below.
\SetLabelAlign{center}{\makebox[\labelwidth][c]{#1}}
\setlist[itemize]{leftmargin=1em, itemsep=0.2em, label={\raisebox{-0.1ex}{\large\textbullet}},
    align=center, labelwidth=0.9em, labelsep=0.4em}
\setlist[enumerate]{leftmargin=1.5em, label=\textbf{\arabic*.}, itemsep=0.25em}
\usepackage{parskip}
% --- CONTROL SECTION FONT SIZE ---
\usepackage{titlesec}
\titleformat{\section}{\large\bfseries}{}{0em}{}
\titleformat{\subsection}{\normalsize\bfseries}{}{0em}{}
\titlespacing*{\subsection}{0pt}{0.75em}{0.25em}
\titlespacing*{\section}{0pt}{1em}{0.5em}
% --- FONT SETUP ---
\usepackage{helvet}
\renewcommand{\familydefault}{\sfdefault}
% --- CUSTOM COMMAND FOR TITLE ---
\newcommand{\makeRecipeTitle}[1]{
    {\noindent \Large \textbf{#1}}
    \par \vspace{-0.75em}
}
% --- CUSTOM COMMAND FOR SERVINGS ---
% Usage: \recipeServings{4--5} or \recipeServings{6 (one burrito each)}
% Place directly after the first \makeRecipeTitle on page 1.
\newcommand{\recipeServings}[1]{
    {\noindent \itshape Serves #1 \par}
    \vspace{-0.75em}
}
% --- PREP ICONS ---
% Usage: \item[\veg] Red Bell Pepper – 1, sliced
%        \item[\measure] Fish Sauce – 1 1/2 tbsp (22ml)
% Replace the bullet on lines to handle before cooking starts: produce to cut
% (carrot), and spices, seasonings, sauces, and pastes with a measured amount
% (measuring spoon, drawn with TikZ). "To taste" lines get neither.
\usepackage{fontawesome5}
\usepackage{tikz}
\newcommand{\veg}{\raisebox{0.25ex}{\small\faCarrot}}
\newcommand{\measure}{\tikz[x=0.9em, y=0.9em,
        baseline={([yshift=-0.36em]current bounding box.center)}]{
    \begin{scope}[rotate=45]
        % Bowl with a rim highlight, and a flat handle with a hanging hole.
        \fill (0,0) ellipse [x radius=0.33, y radius=0.29];
        \draw[white, line width=0.035em, line cap=round]
            (-0.1,0.173) arc[start angle=120, end angle=210, radius=0.2];
        \fill (0.2,-0.07) -- (0.86,-0.105) arc[start angle=-90, end angle=90, radius=0.105]
            -- (0.2,0.07) -- cycle;
        \fill[white] (0.86,0) circle (0.04);
    \end{scope}
}}
% --- CUSTOM COMMAND FOR SOURCE CREDIT ---
% Usage: \recipeSource{Creator -- Video Title}{https://...}
% Place just before \end{document}; pins the credit to the bottom of the last page.
\usepackage[hidelinks]{hyperref}
\urlstyle{same}
\newcommand{\recipeSource}[2]{
    \vfill
    {\noindent \footnotesize \textbf{Source:} #1 \\ \url{#2}}
}


\begin{document}

% ================= PAGE 1 =================
\makeRecipeTitle{Chicken Chanko Nabe (5-Day Meal Prep)}
\recipeServings{5 (meal-prep containers)}

\section*{Ingredients}

\textbf{The Meatballs \& Chicken}
\begin{itemize}[itemsep=0.1em, label={\raisebox{-0.1ex}{\large\textbullet}}]
      \item Ground Chicken – 1 lb (454g), 97\% lean
      \item Egg – 1 large
      \item[\cut] Ginger – 1 tbsp, grated
      \item[\cut] Green Onions – 3, finely chopped
      \item[\measure] Cornstarch – 1 tbsp
      \item[\measure] Soy Sauce – 1 tbsp (15ml)
      \item[\measure] Neutral Oil – 1 tbsp (keeps the lean meatballs moist)
      \item[\cut] Boneless Skinless Chicken Thighs – 3 (about 12 oz / 340g), cut into 1-inch pieces
\end{itemize}

\textbf{The Broth}
\begin{itemize}[itemsep=0.1em, label={\raisebox{-0.1ex}{\large\textbullet}}]
      \item[\measure] Water – 10 cups (2.4L)
      \item[\measure] Dashi Powder – 1 tbsp (Hondashi recommended)
      \item[\measure] Soy Sauce – 1/3 cup (80ml), plus more to taste
      \item[\measure] Mirin – 1/4 cup (60ml)
      \item[\measure] Sake – 1/4 cup (60ml), or dry sherry
      \item[\cut] Garlic Cloves – 3, minced
      \item[\measure] White Miso – 2 tbsp
\end{itemize}

\newpage
\makeRecipeTitle{Chicken Chanko Nabe (5-Day Meal Prep)}
\section*{Ingredients (Continued)}

\textbf{Vegetables, Tofu \& Udon}
\begin{itemize}[itemsep=0.1em, label={\raisebox{-0.1ex}{\large\textbullet}}]
      \item[\cut] Daikon – 8 oz (225g), peeled and cut into 1/4-inch (6mm) half-moons
      \item[\cut] Carrots – 2, peeled and sliced into 1/4-inch (6mm) coins
      \item[\cut] Napa Cabbage – 1.5 lbs (680g), cut into 2-inch pieces, stems and leaves kept apart
      \item[\cut] Shiitake Mushrooms – 8, stems removed, caps left whole (or 5 oz pre-sliced)
      \item[\cut] Leeks – 2, white and light green parts only, sliced 1/2-inch (1cm) thick on the diagonal
      \item[\cut] Firm Tofu – 14 oz (400g), cut into 1-inch cubes
      \item Dried Udon – 1 pack, 14 oz (about 2 oz / 55g per bowl)
      \item Shichimi Togarashi or Sesame Oil – to taste (optional)
\end{itemize}

\section*{Notes}

\begin{itemize}
      \item \textbf{Storage:} Keeps 3--4 days in the refrigerator, which is why the day 4 and 5 containers go in the freezer. Never store the udon in the broth; it turns mushy.
\end{itemize}

\newpage
\makeRecipeTitle{Chicken Chanko Nabe (5-Day Meal Prep)}
\section*{Instructions}

\begin{enumerate}

      \item \textbf{Prep Everything} \\
            Grate the ginger, chop the green onions, and mince the garlic. Cut the vegetables,
            chicken, and tofu as listed, keeping the cabbage stems apart from the leaves. Rinse
            the sliced leeks in a bowl of water to wash out grit, then drain. Measure out
            everything marked with a spoon.

      \item \textbf{Make the Meatball Mix} \\
            In a bowl, combine the ground chicken, egg, ginger, green onions, cornstarch, soy
            sauce, and oil. Mix by hand until sticky, then cover and refrigerate until step 5
            (chilled mix is easier to scoop).

      \item \textbf{Build the Broth} \\
            In a large pot (8 qt is ideal, 6 qt at minimum), combine the water, dashi powder, soy
            sauce, mirin, sake, and garlic. Bring to a boil over high heat, then lower to
            medium-low for a gentle simmer.

      \item \textbf{Cook the Root Vegetables} \\
            Add the daikon and carrots. Simmer until fork-tender but still firm (10 minutes).

      \item \textbf{Add the Chicken} \\
            Add the thighs and simmer (5 minutes). Spoon in the meatball mix, about 2 tbsp each
            (around 16), and simmer until they reach 165$^\circ$F (10 minutes).

      \item \textbf{Finish the Pot} \\
            Add the cabbage stems, shiitake, leeks, and tofu and simmer (3 minutes). Push the
            cabbage leaves under the broth and simmer (2 minutes). Take the pot off the heat. In a
            small bowl, whisk the miso with a ladle of broth until smooth, then stir it into the
            pot. Taste; add soy sauce 1 tsp at a time if needed.

      \item \textbf{Cool and Portion} \\
            Let the pot cool (20 minutes). Divide the solids evenly between 5 containers (1 qt /
            1L each) and ladle the broth over. Refrigerate 3; freeze 2 for days 4--5.

      \item \textbf{Serve} \\
            If the container is frozen, thaw it in the fridge overnight. Boil 2 oz (55g) of udon
            in a separate pot of water (8 minutes), drain, and put it in a bowl. Meanwhile, heat
            the chanko in a small pot until it just simmers (a hard boil dulls the miso), or
            microwave it until steaming (3--4 minutes, stirring halfway). Pour it over the
            noodles and top with shichimi togarashi or sesame oil.

\end{enumerate}

\end{document}
